Vulnerability prioritization methods are routinely validated against CISA's
Known Exploited Vulnerabilities (KEV) catalogue, as if KEV listed every
exploited vulnerability. We measure what that assumption costs, using NIST's
proposed Likely Exploited Vulnerabilities (LEV) metric as a case study. LEV
compounds a vulnerability's history of Exploit Prediction Scoring System
(EPSS) scores into a probability that it has already been exploited; we found
no published evaluation of it. We reimplement LEV, reproduce the
white paper's averages for vulnerabilities (CVEs) published since 2021 to
within 7\%, and score it on 105,250 CVEs. Against KEV, LEV over-predicts by a
factor of 17.3 but still ranks CVEs well. Part of the gap lies in its
inputs: EPSS model releases rescale scores, which accounts for
27--62\% of LEV's mass, and EPSS scores for CVEs that no list records rise
rather than decay, so LEV does not saturate as designed. We then audit public exploitation lists
for dependence and apply capture-recapture estimation. Among 54,612 CVEs,
KEV records at most 32\% of the exploited CVEs that three evidence-citing
lists observe. For CVEs with LEV of at least 0.9, where the estimate is
identified, LEV is about twice as confident as the estimated exploitation
rate, not the 4.5 times that KEV alone suggests. A monotone correction
restores calibration a year ahead against the broader ground truth but must be
refitted as CVEs age. We release code, data where licences allow, and
per-CVE truth bounds so that other methods can be re-scored against more than
KEV.
